\documentclass{article}
\usepackage{amsmath}
\usepackage{graphicx}
\usepackage{hyperref}

\title{A Survey of Vector Databases}
\author{Review Team}
\date{2026}

\begin{document}
\maketitle

\begin{abstract}
We survey 15+ vector database systems across four axes: index type, distance metric, hybrid search support, and operational maturity.
\end{abstract}

\section{Introduction}
Vector databases have emerged as a critical component of modern retrieval-augmented AI pipelines. This survey organises the design space.

\section{Index Types}

\subsection{Tree-based}
Systems like Annoy use random-projection trees.

\subsection{Graph-based}
HNSW has become the de-facto standard. Systems based on HNSW include Faiss, Milvus, and Weaviate.

\subsection{Quantisation}
Product quantisation trades recall for memory. IVF-PQ combines inverted files with PQ for scale.

\section{Distance Metrics}
Most systems support cosine, Euclidean (L2), and inner-product. Some also support Hamming, Jaccard, and custom metrics.

\section{Hybrid Search}
Hybrid search combines vector similarity with lexical (BM25) or structured filters. Approaches vary in how they fuse scores.

\section{Operational Maturity}

\begin{itemize}
  \item Replication: most mature systems support synchronous and asynchronous replication.
  \item Consistency: eventual consistency is common; strict consistency requires careful design.
  \item Scale: petabyte-scale deployments exist but are rare outside hyperscalers.
\end{itemize}

\section{Conclusion}
The field is consolidating around HNSW-based approaches. Differentiation now comes from hybrid-search capabilities and operational tooling.

\section*{Related Surveys}
See \cite{survey2022,survey2023,survey2024} for prior reviews.

\bibliographystyle{plain}
\bibliography{refs}

\end{document}
